% TOP_PROBLEM: {"id": 30007279, "title": "The seven-eighths zero-free half-plane", "category_id": 1, "category": {"id": 1, "name": "number_theory", "display_name": "Number Theory", "description": "Properties of integers, prime numbers, Diophantine equations.", "slug": "number-theory", "order_index": 1, "created_at": "2026-07-31T15:26:25.670Z"}, "status": "open"}
% FIRST_POSTED: 2026-10-06
% ENTRY_KIND: statement_only
% !TeX program = lualatex
\documentclass[11pt]{article}
\usepackage[margin=25mm]{geometry}
\usepackage{fontspec}
\setmainfont{Latin Modern Roman}
\usepackage{amsmath,amssymb,mathtools}
\usepackage{unicode-math}
\setmathfont{Latin Modern Math}
\usepackage{xurl,hyperref}
\hypersetup{colorlinks=true,urlcolor=blue}
\setcounter{secnumdepth}{0}
\setlength{\parindent}{0pt}
\setlength{\parskip}{5pt}
\setlength{\emergencystretch}{3em}
\newcommand{\R}{\mathbb R}
\newcommand{\N}{\mathbb N}
\newcommand{\Z}{\mathbb Z}
\newcommand{\Q}{\mathbb Q}
\newcommand{\C}{\mathbb C}
\newcommand{\D}{\mathbb D}
\newcommand{\F}{\mathbb F}
\newcommand{\sref}[2]{\hyperlink{source:#1:#2}{[S#2]}}
\newcommand{\cataloguescope}[1]{#1}
\begin{document}
\section*{The seven-eighths zero-free half-plane}\label{30007279}
\subsection{Definitions and mathematical statement}
\paragraph{Statement recorded in the supplied source.}
Do all Dirichlet L-functions, including zeta, and all finite-order Hecke L-functions over Q(sqrt(-3)) have no zeros in $Re(s)>7/8$, allowing a pole at one for the principal character?
\subsection{Short English statement}
Do all Dirichlet L-functions, including zeta, and all finite-order Hecke L-functions over Q(sqrt(-3)) have no zeros in $Re(s)>7/8$, allowing a pole at one for the principal character?
\subsection{Sources}
\begin{itemize}
\item[S1] Original question source or mathematical context.
\url{http://www.openproblemgarden.org/op/the_riemann_hypothesis}.
\item[S2] Original question source or mathematical context.
\url{http://en.wikipedia.org/wiki/Riemann_hypothesis}.
\item[S3] Original question source or mathematical context.
\url{http://en.wikipedia.org/wiki/Selberg_class}.
\item[S4] Original question source or mathematical context.
\url{http://www.claymath.org/news/sarnak.php}.
\item[S5] Original question source or mathematical context.
\url{http://www.claymath.org/millennium/Riemann_Hypothesis/}.
\item[S6] Original question source or mathematical context.
\url{http://empslocal.ex.ac.uk/people/staff/mrwatkin/zeta/RHproofs.htm}.
\item[S7] Original question source or mathematical context.
\url{http://vixra.org/pdf/1210.0176v7.pdf}.
\item[S8] Original question source or mathematical context.
\url{http://arxiv.org/pdf/1210.1517v10.pdf}.
\item[S9] Original question source or mathematical context.
\url{https://doi.org/10.4171/owr/2020/17}.
\item[S10] Original question source or mathematical context.
\url{https://dlmf.nist.gov/25.15}.
\item[S11] Original question source or mathematical context.
\url{https://arxiv.org/abs/2608.15773}.
\item[S12] Original question source or mathematical context.
\url{https://www.preprints.org/manuscript/202506.0481}.
\item[S13] OpenAI, manuscript citation page. Source of the formulation motivating this local problem.
\url{https://github.com/openai/math/blob/adc7f1241b42e322a6451854ab7e4b4c146bf78a/preprints/The-Quasi-Riemann-Hypothesis-September-30-2026/README.md}.
\end{itemize}
\end{document}
